\section{Summary of Results}
\label{sec:results}

Figure~\ref{fig:ladder} puts the whole evolution on one axis.

\begin{figure}[h]
  \centering
  \includegraphics[width=\textwidth]{fig_gap_ladder}
  \caption{Closing the gap (run W): scalar C++ throughput over each OCaml
  variant's. Every build keeps its polls (on back edges, or at loop exits with
  counted loops). The SSE2 bar adds one change to the last ladder step.
  Abbreviations: s.r., strength reduction; i.p., interior pointers;
  inv.\ str., invariant-stride induction variables.}
  \label{fig:ladder}
\end{figure}

\subsection{Every Final Configuration}
\label{sec:unified}

\begin{table}[h]
  \centering
  \small
  \begin{tabularx}{\textwidth}{@{}Xrrr@{}}
    \toprule
    configuration (same OCaml source unless noted) & tokens/s & vs scalar C++ & vs vectorised C++ \\
    \midrule
    C++ vectorised (\code{cpp_unroll}) & @T{cpp_unroll}@ & @G{cpp_unroll}@ & $1.00\times$ \\
    C++ scalar (\code{cpp_novec}) & @T{cpp_novec}@ & $1.00\times$ & @GU{cpp_novec}@ \\
    Rust (\code{rust_unroll}, bounds-checked) & @T{rust_unroll}@ & @G{rust_unroll}@ & @GU{rust_unroll}@ \\
    \midrule
    stock OCaml 5.3.0, first port (\code{ocaml_safe_plain}) & @T{ocaml_safe_plain}@ & @G{ocaml_safe_plain}@ & @GU{ocaml_safe_plain}@ \\
    \quad \code{-unsafe}, unrolled, hand s.r.\ (\code{ocaml_unsafe_sr}) & @T{ocaml_unsafe_sr}@ & @G{ocaml_unsafe_sr}@ & @GU{ocaml_unsafe_sr}@ \\
    \quad \code{-unsafe}, unrolled, ported passes + counted loops (\code{ocaml_pr_counted}) & @T{ocaml_pr_counted}@ & @G{ocaml_pr_counted}@ & @GU{ocaml_pr_counted}@ \\
    \midrule
    OxCaml 5.4.0-ox7 release, first port (\code{oxcaml_safe_plain}) & @T{oxcaml_safe_plain}@ & @G{oxcaml_safe_plain}@ & @GU{oxcaml_safe_plain}@ \\
    \quad \code{-unsafe}, unrolled (\code{oxcaml_unsafe_unroll}) & @T{oxcaml_unsafe_unroll}@ & @G{oxcaml_unsafe_unroll}@ & @GU{oxcaml_unsafe_unroll}@ \\
    \quad + hand s.r.\ (\code{oxcaml_unsafe_sr}) & @T{oxcaml_unsafe_sr}@ & @G{oxcaml_unsafe_sr}@ & @GU{oxcaml_unsafe_sr}@ \\
    \midrule
    PR compiler, \code{-unsafe}, unrolled, passes off (\code{pr_off}) & @T{pr_off}@ & @G{pr_off}@ & @GU{pr_off}@ \\
    \quad all passes + \code{-cmm-counted-loops} (\code{pr_counted}) & \textbf{@T{pr_counted}@} & \textbf{@G{pr_counted}@} & \textbf{@GU{pr_counted}@} \\
    \quad + SSE2 kernel (\code{pr_simd_counted}) & @T{pr_simd_counted}@ & @G{pr_simd_counted}@ & @GU{pr_simd_counted}@ \\
    \midrule
    bounds checks kept: release, unrolled (\code{oxcaml_safe_unroll}) & @T{oxcaml_safe_unroll}@ & @G{oxcaml_safe_unroll}@ & @GU{oxcaml_safe_unroll}@ \\
    \quad PR compiler, all passes + counted loops (\code{pr_safe_counted}) & @T{pr_safe_counted}@ & @G{pr_safe_counted}@ & @GU{pr_safe_counted}@ \\
    \bottomrule
  \end{tabularx}
  \caption{The core configurations in run W (generation mode, median of 21
  interleaved rounds; ratios are C++ throughput over the variant's; s.r.,
  strength reduction). All rows but the first ports and the last two use
  \code{-unsafe}; the headline is the scalar \code{pr_counted} row, and the
  SSE2 row adds an explicitly vectorised kernel. Close builds can be ordered differently here than by
  the paired comparisons, which are the reliable ones. Scoring mode agrees on
  every difference larger than the resolution.}
  \label{tab:final}
\end{table}

Beyond Table~\ref{tab:final}, paired comparisons settle three questions:

\begin{itemize}
  \item \textbf{Counted loops carry the gain:} @P{pr_all,pr_counted}@
        over the four passes with a poll on every back edge, and with the
        passes @P{pr_off,pr_counted}@ over passes off
        (@Pp{pr_off,pr_counted}@ in scoring mode), and
        @P{oxcaml_unsafe_unroll,pr_counted}@ over the release; the price is
        pause latency (Section~\ref{sec:counted}).
  \item \textbf{The gain needs the pointer passes:} within counted mode they
        are worth @P{pr_counted_noptr,pr_counted}@; moving polls to loop exits
        matters because it makes the pointers legal.
  \item \textbf{Counted loops are as fast as poll-free code} within the run's
        resolution, @D{pr_counted,pr_nopoll}@ to no polls at all
        (Section~\ref{sec:last2}); unchecked loops and the other builds that
        give up liveness gain no more than about the resolution
        (Table~\ref{tab:liveness}).
\end{itemize}

\paragraph{With bounds checks kept.} Every OCaml figure above except the
first ports uses \code{-unsafe}; Rust keeps its bounds checks and is
@G{rust_unroll}@ behind C++. With checks kept, the pull request is
\emph{slower} than the release: @P{oxcaml_safe_unroll,pr_safe_counted}@ against
bounds-checked OxCaml, @G{pr_safe_counted}@ behind C++, and
@P{pr_counted,pr_safe_counted}@ against the counted build without checks.
Interior pointers still form, but each check needs the row's index as well as
its pointer: with eight values live instead of four, the allocator spills and
reloads the four row pointers every iteration. I did not build a fix.

\subsection{Does the Gap Survive a Wider Model?}
\label{sec:width}
microgpt's weights fit in the L1 data cache (Section~\ref{sec:model}). A
sweep (a separate run, 21 interleaved rounds) widened the model, with no code
change, to one layer, four heads and widths of 16 to 256, taking the weights
through L2 (116\,KB--1.6\,MB) into L3 (6.4\,MB); all variants agree on every
model's output. The wider models have random weights, which generate names
about twice as long, so attention does more work per token; at width 16 that
alone moves the counted build from $1.24\times$ to $1.42\times$ behind scalar
C++. Widths are therefore compared among random models
(Figure~\ref{fig:width}).

\begin{figure}[h]
  \centering
  \includegraphics[width=\textwidth]{fig_width}
  \caption{The gap to scalar C++ as the model widens (width-sweep run, one
  layer, 21 rounds, median tokens/s). Lines are models with random weights;
  the open markers are the trained width-16 model. The patched compiler with
  counted loops polls at loop exits.}
  \label{fig:width}
\end{figure}

With counted loops, OxCaml goes from $1.42\times$ behind scalar C++ at width
16 to $1.32\times$ at 32, $1.26\times$ at 64, $1.16\times$ at 128 and
$1.10\times$ at 256 (C++ won at least 20 of 21 rounds at every width); the
5.2.0+ox release build stays $1.66$--$1.93\times$ behind. At widths 32--256
counted loops gain $+46$ to $+51\%$ over that build and $+27$ to $+31\%$ over
the four passes alone (paired medians, at least 20/21 rounds each).

Why does the gap narrow? If memory is the cause, the ratio should narrow only
when the weights stream from a slower level; if the shape is the cause
(longer rows spreading each row's setup over more multiply--adds), it should
narrow with matrix size even in cache. A separate micro-benchmark of the
four-row kernel, with the weights reused (in cache) or cycled through 8\,MB
(L3) or 64\,MB (main memory), decides it. The OCaml/C++ ratio falls from
$1.19\times$ at $16\times16$ to $1.07$--$1.14\times$ at $256\times256$ and
larger in cache, and no further from L3; only from main memory, which no
model reaches, do the two converge. The narrowing therefore follows the loop
shape: longer rows amortise the setup that Section~\ref{sec:residual} places
outside the inner loop. GCC's vectoriser stays worth 8--10\% at every width,
and Rust within 1--7\% of scalar C++.

\subsection{In Words}
Source-level work took OxCaml from @G{oxcaml_safe_plain}@ to
@G{oxcaml_unsafe_unroll}@ behind scalar C++. The compiler work, on unchanged
source, closed more than half of the remaining gap, to @G{pr_counted}@ with
counted loops (@GU{pr_counted}@ behind vectorised C++).
Figure~\ref{fig:matrix} shows every core build of run W.

\begin{figure}[t]
  \centering
  \includegraphics[width=\textwidth]{fig_matrix}
  \caption{The core builds of run W (median of 21 interleaved rounds; Python,
  @T{python}@ tokens/s, is omitted). \texttt{oxcaml\_*} builds use the
  5.4.0-ox7 release and \texttt{pr\_*} the pull-request compiler
  (Section~\ref{sec:method}).}
  \label{fig:matrix}
\end{figure}
